
\begingroup
\shorthandoff{<>}% Babel español trata de otro modo ciertos delimitadores de Xy-pic.
\begin{theoreme}
\label{XIII.2.4}
Sean
\marginpar{380}
$f\colon X\to S$ un $S$-esquema, $D$ un divisor sobre~$X$
de cruzamientos normales relativamente a~$S$ \eqref{XIII.2.1},
$Y=\Supp D$, $U=X-Y$, e $i\colon U\to X$ la inmersión canónica. Sea $F$
un haz de conjuntos (\resp de grupos) sobre~$U$ que satisfaga una de las
condiciones siguientes:
\begin{enumerate}
\item[a)] $F$ es, localmente para la topología étale \emph{sobre}
$X$ y sobre~$S$, la imagen inversa de un haz de conjuntos (\resp de un
haz de grupos constructible) sobre~$S$.
\item[b)] $F$ es localmente constante y constructible sobre~$U$, y está
moderadamente ramificado sobre~$X$ relativamente a~$S$.
\end{enumerate}
Entonces se tienen las conclusiones siguientes:
\begin{enumerate}
\item[1)]$(F,i)$ es cohomológicamente propio relativamente a~$S$
en dimensión $\leq 0$ (\resp para todo morfismo $h\colon S'\to S$, si
$i'\colon U'\to X'$ es la imagen inversa de~$i$ sobre~$S'$, si $F'=F|U'$
y si $k=h_{(X)}$, el morfismo canónico
$$
\Psi\colon k^*(\R^1_{\tame}i_*F)\to\R^1_{\tame}i_*'F'
$$
es un isomorfismo).

Si $F$ es un haz de grupos primo con las características
residuales de~$S$
\ifthenelse{\boolean{orig}}
{(\Ref{XIII.2.3}.b))}
{(\Ref{XIII.2.3}.b)}
(moderadamente
ramificado sobre~$X$ relativamente a~$S$), entonces $(F,i)$ es
cohomológicamente propio relativamente a~$S$ en dimensión $\leq 1$.
\item[2)]Si $F$ es un haz de conjuntos (\resp de grupos)
constructible, $i_*F$ (\resp $\R^1_{\tame}i_*F$) es constructible.
\end{enumerate}
\end{theoreme}

\subsubsection*{Demostración}
Para todo $S$-esquema $S'$, consideremos el diagrama siguiente,
en el que todos los cuadrados son cartesianos:
$$
\xymatrix{ U \ar[dd]_-g \ar[dr]^-i & & & U'\ar[lll] \ar[dd]_{\hbox{\raise10mm\hbox{$g'$}}}
\ar[dr]^-{i'}& \\ & X \ar[dl]_{f\mkern-8mu} & & & X'\ar[lll]_{k\hspace*{1cm}} \ar[dl]_{f'\mkern-10mu} \\ S
& & & S'\ar[lll]_-h & & }
$$
Como
\marginpar{381}
la cuestión es local sobre~$X$ para la topología étale, puede suponerse
que $D$ es un divisor de cruzamientos normales estrictos relativamente
a~$S$ \eqref{XIII.2.1}; además, tras restringir si es preciso $X$ a un
entorno de~$Y$, puede suponerse que $X$ es liso sobre~$S$.

\subsubsection*{Demostración de~\Ref{XIII.2.4}~\textup{1)}}
\setcounter{subsection}{4}% Los párrafos numerados pertenecen al Teorema XIII.2.4.
\setcounter{subsubsection}{0}% Comenzar visiblemente por XIII.2.4.1.
\subsubsection{}
\label{XIII.2.4.1}
\emph{Caso de un haz de conjuntos que satisface} a). Puede suponerse que
$F=g^*G$, donde $G$ es un haz sobre~$S$. Se deduce entonces de
(SGA~4 XVI~3.2) que el morfismo canónico
\begin{equation*}
\label{eq:XIII.2.4.1.1}
\tag{\thesubsubsection.1} f^*G\to i_*F
\end{equation*}
es un isomorfismo. Para todo $S$-esquema $h\colon S'\to S$, se tiene
asimismo un isomorfismo $f'^*G'\to i_*'F'$; por consiguiente, el morfismo
canónico
$$
\varphi\colon k^*(i_*F)\to i_*'F'
$$
se identifica con el isomorfismo natural
$$
k^*f^*G\simeq f'^*G'
$$

\subsubsection{}
\label{XIII.2.4.2}
\emph{Caso de un haz de conjuntos que satisface} b). Hay que demostrar
que $\varphi$ es un isomorfismo, y para ello basta comprobarlo en cada
punto geométrico $\overline{x}'$ de~$X'$. Sea $\overline{S}$
(\resp $\overline{X}$, \resp $\overline{S'}$, \resp $\overline{X'}$)
la localización estricta de~$S$ (\resp de~$X$, \resp de~$S'$,
\resp de~$X'$) en $\overline{x}'$, y pongamos
$\overline{U}=U_{(\overline{X})}$,
$\overline{U'}=U_{(\overline{X'})}$, etc. El morfismo
\ifthenelse{\boolean{orig}}
{$\varphi_{\overline{x}}$}
{$\varphi_{\overline{x}'}$}
se identifica con el morfismo canónico
$$
\overline{\varphi}\colon \H^0(\overline{U},\overline{F})\to
\H^0(\overline{U'}, \overline{F'})
$$
Puede encontrarse un recubrimiento principal $V$ de~$\overline{U}$, del
tipo que figura en~\Ref{XIII.5.4}, tal que la imagen inversa de
$\overline{F}$ sobre~$V$ sea un haz constante de valor $C$. Si
$\Pi$ es el grupo de Galois de~$V$ sobre~$\overline{U}$, $\Pi$
actúa sobre~$\overline{F}|V$, y se tiene
\begin{equation*}
\label{eq:XIII.2.4.2.1}
\tag{\thesubsubsection.1}
\ifthenelse{\boolean{orig}}
{\H^0(U,F))\simeq \H^0(V,C_V)^{\Pi}}
{\H^0(\overline{U},\overline{F})\simeq \H^0(V,C_V)^{\Pi}}
\end{equation*}
\label{indnot:mh}\oldindexnot{$\H^0(V,C_V)^{\Pi}$|hyperpage}%
donde el segundo miembro designa el conjunto de los elementos
de~$\H^0(V,C_V)$
\marginpar{382}
invariantes bajo la acción de~$\Pi$. Como
$V'=V\times_{\overline{U}}\overline{U'}$ es un recubrimiento
principal de~$\overline{U'}$, de grupo de Galois $\Pi'\simeq \Pi$, se ve que
el morfismo $\overline{\varphi}$ se obtiene, tomando los invariantes
bajo la acción de~$\Pi$, a partir del morfismo canónico
$$
\H^0(V,C_V)\to \H^0(V',C_{V'}).
$$
Como $V$ y $V'$ son conexos \eqref{XIII.5.4}, este morfismo, y por tanto
también $\overline{\varphi}$, es un isomorfismo.

Obsérvese que, si además $F$ es un haz de grupos y $P$ es un torsor
sobre~$\overline{U}$ bajo $\overline{F}$, moderadamente ramificado
relativamente a~$D$, se deduce de la demostración precedente y
de~\Ref{XIII.2.2} que
$(\leftexp{P}{\overline{F}},\overline{i})$ es cohomológicamente
propio relativamente a~$S$ en dimensión $\leq 0$.

\subsubsection{}
\label{XIII.2.4.3}
\emph{Caso de un haz de grupos.} Para demostrar que $\Psi$ es un
isomorfismo, basta probar que, para todo punto geométrico
$\overline{y}'$ de~$Y'$, el morfismo
$$
\Psi_{\overline{y}'}\colon
(k^*(\R^1_{\tame}i_*F))_{\overline{y}'}\to (\R^1_{\tame}
i_*'F')_{\overline{y}'}
$$
es un isomorfismo. Ahora bien, según~\Ref{XIII.2.1.2},
$\Psi_{\overline{y}'}$ se identifica con el morfismo canónico
$$
\overline{\Psi}\colon
\H^1_{\tame}(\overline{U},\overline{F})\to
\H^1_{\tame}(\overline{U'},\overline{F'}).
$$
Sean $\widetilde{U}$ el recubrimiento universal moderadamente
ramificado de~$\overline{U}$ \eqref{XIII.2.1.3} y $\widetilde{F}$
la imagen inversa de~$\overline{F}$ sobre~$\widetilde{U}$. De
\Ref{XIII.5.7} en el caso a) y de~\Ref{XIII.5.5} en el caso b) resulta
que $\H^1_{\tame}(\overline{U},\overline{F})$ se identifica con el
subconjunto de~$\H^1(\overline{U},\overline{F})$ formado por las clases
de torsores bajo~$\overline{F}$ cuya imagen inversa sobre~$\widetilde{U}$
es trivial. Por otra parte, un razonamiento
\ifthenelse{\boolean{orig}}
{clásico}
{clásico (\cf \Ref{IX.5}, p\ptbl\pageref{page-300})}
muestra que el conjunto de los elementos
de~$\H^1(\overline{U},\overline{F})$ cuya imagen inversa
sobre~$\widetilde{U}$ es trivial se identifica con
$\H^1(\pi_1^{\tame}(\overline{U}),\H^0(\widetilde{U},\widetilde{F})).$
Se obtiene así un isomorfismo canónico
\begin{equation*}
\label{eq:XIII.2.4.3.1} \tag{\thesubsubsection.1} {}
\H^1_{\tame}(\overline{U},\overline{F})\isomto
\H^1(\pi_1^{\tame}(\overline{U}),\H^0(\widetilde{U},\widetilde{F})).
\end{equation*}

Por consiguiente, el morfismo $\overline{\Psi}$ se identifica con el
morfismo canónico
$$
\H^1(\pi_1^{\tame}(\overline{U}),\H^0(\widetilde{U},\widetilde{F}))
\to
\H^1(\pi_1^{\tame}(\overline{U'}),\H^0(\widetilde{U}',\widetilde{F}')).
$$
Demostremos
\marginpar{383}
que este morfismo es un isomorfismo. El morfismo
$\pi_1^{\tame}(\overline{U'})\to\pi_1^{\tame}(\overline{U})$
es un isomorfismo por~\Ref{XIII.5.6}, y lo mismo ocurre con el morfismo
$\H^0(\widetilde{U},\widetilde{F})\to\allowbreak
\H^0(\widetilde{U}',\widetilde{F}')$. En efecto, esto es evidente en el
caso b), pues $\widetilde{F}$ es constante y $\widetilde{U}$ y
$\widetilde{U}'$ son conexos. En el caso a), sea $\overline{G}$ un haz
de grupos constructible sobre~$\overline{S}$ tal que
$\overline{F}=\overline{g}^*\overline{G};$ como los morfismos
$\widetilde{U}\to\overline{S}$ y
$\widetilde{U}'\to\overline{S'}$ son $0$-acíclicos
\eqref{XIII.5.7}, se tiene
$$
\H^0(\widetilde{U},\widetilde{F})\simeq
\H^0(\overline{S},\overline{G}) \simeq
\H^0(\overline{S'},\overline{G'})\simeq
\H^0(\widetilde{U}',\widetilde{F}'),
$$
lo que implica que $\Psi$ es un isomorfismo. La última afirmación
de~\Ref{XIII.2.4}~1) se deduce de lo anterior, teniendo en cuenta
\Ref{XIII.2.3}~b).

\subsubsection*{Demostración de~\Ref{XIII.2.4}~\textup{2)}}

El caso de un haz de conjuntos constructible que satisface a) se deduce
inmediatamente de~\eqref{eq:XIII.2.4.1.1}. Sea $F=g^*G$ un haz de grupos
que satisface a), donde $G$ es un haz constructible; puede suponerse que
$S$ es afín; sea $(S_j)_{j\in J}$
\ifthenelse{\boolean{orig}}
{una familia finita de subesquemas cerrados reducidos de~$S$}
{una familia finita de subesquemas localmente cerrados reducidos de~$S$}
cuya unión recubre~$S$, y tal que la imagen inversa de~$G$ sobre~$S_j$
sea un haz localmente constante. Teniendo en cuenta
\Ref{XIII.2.4}~1), para establecer que $\R^1_{\tame}i_*F$ es
constructible basta comprobarlo después del cambio de base
$S_j\to S$, para cada $j\in J$. Quedamos así reducidos al caso~b),
en el que $F$ es localmente constante.

Supongamos en adelante que $F$ es un haz de conjuntos o de grupos que
satisface b). Como la cuestión es local para la topología étale sobre~$X$,
puede suponerse que $X$ es de presentación finita sobre~$S$ y, por paso
al límite, que $X$ y $S$ son noetherianos.

Sea $D=\sum_{1\leq i\leq r}\divisor f_i$, donde, para cada punto
$x$ de~$\Supp D$, si $I(x)$ es el conjunto de los $i$ tales que
$f_i(x)=0$, el subesquema $V((f_i)_{i\in I(x)})$ es liso sobre~$S$ de
codimensión $\card I(x)$ en~$X$. Sea $\mathcal{P}$ el conjunto de las
partes de~$[ 1,r ]$ y, para cada $I\in\mathcal{P}$, pongamos
$$
X_I=\textstyle\Big(\bigcap_{i\in I}V(f_i)\Big)\cap\Big(\bigcap_{i\not\in I}X_{f_i}\Big).
$$
Sea
\marginpar{384}
$z$ un punto de~$X_I$. Tras restringirnos primero a un entorno étale
de~$z$, pueden encontrarse un abierto $W$ de~$X$ que contiene a~$z$ y un
recubrimiento principal $V$ de~$U\cap W$, moderadamente ramificado
sobre~$W$ relativamente a~$S$, del tipo considerado
en~\Ref{XIII.5.6.1}, tal que la imagen inversa de~$F$ sobre~$V$ sea
un haz constante de valor~$C$. Sea $\pi$ el grupo de Galois del
recubrimiento~$V$. Para todo punto geométrico $\overline{x}$
\ifthenelse{\boolean{orig}}
{de~$X_I$}
{de~$X_I\cap W$}
se tiene entonces, por~\eqref{eq:XIII.2.4.2.1},
$$
\H^0(\overline{U},\overline{F})\simeq \H^0(V,C_V)^{\pi}.
$$
De ello resulta que $i_*F|X_I\cap W$ es localmente constante
(SGA~4 IX~2.13) y, por consiguiente, $i_*F$ es constructible.

Demostremos, por último, que, si $F$ es un haz de grupos localmente
constante, $\R^1_{\tame}i_*F|_{X_I}$ es constructible. Si
$\overline{x}$ es un punto geométrico
\ifthenelse{\boolean{orig}}
{de~$X$}
{de~$X_I$},
se obtuvo en~\eqref{eq:XIII.2.4.3.1} la expresión
$$
\H^1_{\tame}(\overline{U},\overline{F})\isomto
\H^1(\pi^{\tame}_1(\overline{U}),\H^0(\widetilde{U},\widetilde{F})).
$$
Si $p$ es la característica residual de~$\overline{X}$, se tiene,
por~\Ref{XIII.5.6},
$\pi_1^{\tame}(\overline{U})=\prod_{l\neq p }\ZZ_l(1)^{\card I}$.
Denotemos por $\LL$ el conjunto de los números primos que dividen el
orden del grupo finito $\H^0(\widetilde{U},\widetilde{F})$, y sea
$K=\prod_{l\in\LL -\{p\}\cap\LL }\ZZ_l(1)^{\card I}$. De
\cite[I \S5 ex.2]{XIII.4} resulta que se tiene
$$
\H^1_{\tame}(\overline{U},\overline{F})\simeq
\H^1(K,\H^0(\widetilde{U},\widetilde{F})).
$$
Como $K$ está topológicamente generado por un número finito de elementos y
$\H^0(\widetilde{U},\widetilde{F})$ es finito, se deduce en primer lugar
que las fibras del haz $\R^1_{\tame}i_*F|_{X_I}$ son finitas. Por otra
parte, el conjunto $\LL$ no depende del punto $\overline{x}$. Para todo
$q\in\LL$, sea $X_{I,q}$ el cerrado de~$X_I$ definido por la ecuación
$q=0$, y sea
\ifthenelse{\boolean{orig}}
{$X_{I'}$}
{$X_I'$}
el abierto de~$X_I$ complementario de la unión de los $X_{I,q}$. Entonces
$\R^1_{\tame}i_*F|_{X_{I,q}}$ y $\R^1_{\tame}i_*F|_{X_I'}$ son
localmente constantes; en efecto, una flecha de especialización de puntos
geométricos de~$X_{I,q}$ (\resp de
\ifthenelse{\boolean{orig}}
{$X_{I'}$}
{$X_I'$})
induce un isomorfismo sobre los grupos $K$ \eqref{XIII.5.6.1}, y por tanto
también sobre los conjuntos
$\H^1_{\tame}(\overline{U},\overline{F})$; puede aplicarse entonces
SGA~4 IX~2.13.
\endgroup
